%*************************************************
% In this file the first few pages are typeset.
% Make changes accordingly.
%*************************************************

% شماره صفحات با حروف
\pagenumbering{adadi}

%***************************
% Page 0: white paper
%***************************
\thispagestyle{empty}
\begin{center}
	~\vfill
%	\includegraphics[scale=1]{images/besm.jpg}
	~\vfill
\end{center}
\pagebreak

%***************************
% Page 1: Besmelah
%***************************
\thispagestyle{empty}
\begin{center}
	~\vfill
	\includegraphics[scale=1]{images/besm.jpg}
	~\vfill
\end{center}
\pagebreak

%***************************
% Page 2: Title
%***************************
\thispagestyle{empty}
\newgeometry{left=3cm,right=3cm,top=2cm}
\begin{center}
\includegraphics[height=3cm]{images/iut_logo.png}
\vspace{0.4cm}

{\large
	\textbf{دانشگاه صنعتی اصفهان}\\
	دانشکده مهندسی برق و کامپیوتر
}
\vspace{3cm}

{\LARGE
	\textbf{
	 پیش‌بینی‌ و پایش بی‌درنگ داده‌های شبیه‌سازی موتورهای جت CMAPSS ناسا بر اساس چارچوب معماری دوقلوهای دیجیتال چندعاملی}\\
}
\vspace{3cm}

{\large
	\textbf{پایان‌نامه کارشناسی ارشد مهندسی کامپیوتر - نرم‌افزار}\\
}
\vspace{1cm}

{\Large
	\textbf{نوید نصیری}\\
}
\vspace{2.5cm}

{\large
	استاد راهنما\\
}
\vspace{0.5cm}

{\Large
	\textbf{دکتر مسعودرضا هاشمی}\\
}
\vspace{1cm}

{\large
	استاد مشاور\\
}
\vspace{0.5cm}

{\Large
	\textbf{دکتر عارف کریمی افشار}\\
}
\vspace{2.5cm}
\vspace{3.75cm}

{\Large
	\textbf{1400}
}

\end{center}
\restoregeometry
\pagebreak

%***************************
% 4th page: Signatures
%***************************
\thispagestyle{empty}
\newgeometry{left=3cm,right=3cm,top=2cm}
\begin{center}
	\includegraphics[height=3cm]{images/iut_logo.png}
	\vspace{0.4cm}
	
	\textbf{دانشگاه صنعتی اصفهان}\\
	\vspace{0.4cm}
	
	{\large
		دانشکده مهندسی برق و کامپیوتر
	}
	\vspace{1.5cm}
	
	\vfill
	
	{\Large
		پایان‌نامه کارشناسی ارشد رشته مهندسی کامپیوتر -- نرم‌افزار آقای نوید نصیری\\
		\vspace{.3cm}
		تحت عنوان\\
	}
	
	
\end{center}
\vfill
\vspace{2.5cm}

{\large
	\textbf{ پیش‌بینی‌ و پایش بی‌درنگ داده‌های شبیه‌سازی موتورهای جت CMAPSS ناسا بر اساس چارچوب معماری دوقلوهای دیجیتال چندعاملی
	}
}

\vspace*{2cm}
%
%در تاریخ 1404/01/01 توسط کمیته تخصصی زیر مورد بررسی و تصویب نهایی قرار گرفت:\\
\vspace{0.8cm}

%{\normalsize
%	\begin{tabular}{rr}
%		1- استاد راهنمای پایان‌نامه & دکتر محمد محمدی محمدآبادی \\[0.8cm]
%		2- استاد مشاور پایان‌نامه & دکتر محمد محمدی محمدآبادی \\[0.8cm]
%		3- استاد داور & دکتر محمد محمدی محمدآبادی \\[0.8cm]
%		4- استاد داور & دکتر محمد محمدی محمدآبادی \\[0.8cm]
%		سرپرست تحصیلات تکمیلی دانشکده & دکتر محمد محمدی محمدآبادی \\[0.8cm]
%	\end{tabular}
%}
\restoregeometry
\pagebreak

%***************************
% 5th page: Acknowledgment
%***************************
\thispagestyle{empty}
\newgeometry{left=3cm,right=4cm,top=4cm}
\vspace*{1.5cm}

{\large
	\textbf{تشکر و قدردانی}\\
	
	
	سپاس خدایی را که با وجود ناتوانی آفریدگان در گزاردن حقش، باز الطاف محبت بی‌کرانش بر همه می‌تابد.
	
	قدردان راهنمایی‌های ارزنده استاد گرانقدر جناب آقای دکتر مسعودرضا هاشمی و جناب آقای دکتر عارف کریمی افشار هستم که بدون شک پشتیبانی و راهنمایی‌هایشان روشنگر مسیر انجام پژوهش بود.
	
}
\restoregeometry
\pagebreak

%***************************
% 6th page: Rights
%***************************
%\thispagestyle{empty}
%\newgeometry{left=6cm,right=6cm}
%
%begin{spacing}{3}
%	\leavevmode
%	\vfill
%	\parbox{8 cm}{
%		
%		\textbf{\Large کلیه حقوق مادی مترتب بر نتایج مطالعات، ابتکارات و نوآوری‌های ناشی از تحقیق موضوع این پایان‌نامه  متعلق به دانشگاه صنعتی اصفهان است.}
%		
%	}
%	\vfill
%\end{spacing}
%\restoregeometry
%\pagebreak

\thispagestyle{empty}
\newgeometry{left=2.5cm,right=3cm}

\vspace*{3cm}
\textbf{\Large کلیه حقوق مالکیت مادی و معنوی مربوط به اين پايان نامه متعلق به دانشگاه صنعتی اصفهان و پدیدآورندگان است. این حقوق توسط دانشگاه صنعتي اصفهان و بر اساس خط مشی مالکیت فکری این دانشگاه، ارزش‌گذاری و سهم بندي خواهد شد.
\\
			هر گونه بهره برداري از محتوا، نتايج یا اقدام براي تجاري‌سازي دستاوردهاي اين پايان نامه تنها با مجوز کتبی دانشگاه صنعتی اصفهان امکان‌پذیر است.
}

\restoregeometry
\pagebreak

%***************************
% 7th page: Dedication
%***************************
\thispagestyle{empty}
\vspace*{3cm}

{\nastaliq
	\Huge
	\raggedleft
	\textbf{\vspace{2cm}تقدیم به}	
	\\
\hspace{1cm}	
\textbf{\vspace{0.5cm}پدر و مادر عزیزم}
\\
\hspace{4cm}
	مستحکم‌ترین پشتوانه و حامیانم که وجودشان منشأ آرامش و انگیزه است...\vspace{1cm}
\\ 
\hspace{1cm}	
\textbf{\vspace{0.5cm}و همسرم}
\\
\hspace{4cm}
	که مشوق و همراه همیشگی من است...

	
}
\pagebreak

%***************************
% 8th page: Table of contents
%***************************

\titleformat{\chapter}[display]
	{\normalfont\LARGE\bfseries\centering}{\chaptertitlename ~ \tartibi{chapter}}{20pt}{\LARGE}
\newgeometry{left=2.5cm,right=3cm,top=3cm,bottom=2.5cm,includehead=false,headsep=1cm,footnotesep=.5cm}
\baselineskip=.7cm


\addtocontents{toc}{\textbf{\underline{عنوان}}}
\addtocontents{toc}{\hfill\textbf{\underline{صفحه}}\par}
\addcontentsline{toc}{section}{فهرست مطالب}
\tableofcontents
\pagebreak



%\begin{comment}
\renewcommand\listfigurename{فهرست شکل‌ها}
\addcontentsline{toc}{section}{فهرست شکل‌ها}
\listoffigures
\pagebreak
%\end{comment}

%\begin{comment}
\renewcommand\listfigurename{فهرست جداول}
\addcontentsline{toc}{section}{فهرست جداول}
\listoftables
\pagebreak
%\end{comment}


% change the font and margins of a chapter title
\titlespacing*{\chapter}{0pt}{3.5cm}{6cm}
\titleformat{\chapter}[display]
	{\normalfont\LARGE\bfseries\raggedright}{\chaptertitlename ~ \tartibi{chapter}}{20pt}{\LARGE}

% No page numbers on the first page of a chapter
\assignpagestyle{\chapter}{empty}
